\documentclass[twocolumn]{aastex631}

\shorttitle{False Positive Characterization of TESS Candidates}
\shortauthors{Mehta}

\title{Multi-Tiered Vetting and False Positive Identification of Four TESS Candidate Systems: TIC 258285711, TIC 319431206, TIC 420914536, and TIC 432121978}

\author[0009-0005-4882-9388]{Bhavish Mehta}
\affiliation{Independent Researcher, Ahmedabad, India}

\begin{abstract}
We report the multi-tiered vetting of four high-priority transiting exoplanet candidates identified in TESS Full Frame Image (FFI) data: TIC 258285711, TIC 319431206, TIC 420914536, and TIC 432121978. Using a diagnostic framework incorporating statistical modeling, kinematic transit duration constraints, Lomb-Scargle out-of-transit stellar variability checks (SWEET test), odd-even transit depth mismatches, and trapezoid vs. triangle U/V shape profiling, we characterize all four systems. We find that all four targets are astrophysical false positives. Specifically, TIC 258285711 and TIC 319431206 are retired due to severe orbital duration violations, with TIC 258285711 showing a phase-locked stellar rotation period. The two short-period candidates, TIC 420914536 and TIC 432121978, which initially passed basic kinematic checks, are retired as grazing eclipsing binaries based on severe odd-even depth mismatches and V-shaped profiles confirmed via multi-model $\chi^2$ fitting.
\end{abstract}

\section{Introduction}
The Transiting Exoplanet Survey Satellite (TESS; Ricker et al. 2015) has detected thousands of planet candidates. However, the large $21''$ pixel scale of the TESS cameras often results in transit signals contaminated by nearby astrophysical false positives, such as blended eclipsing binaries (BEBs) or stellar variability. Automated, multi-tiered vetting pipelines are essential to optimize follow-up resources by efficiently filtering out these false alarms. This Note details the systematic vetting and retirement of four candidates using a combination of orbital diagnostics, odd-even depth checks, and U/V shape profiling.

\section{Methodology}
We performed a multi-tiered vetting analysis on the four targets using TESS FFI light curves stitched across available sectors. Our pipeline evaluated the candidates across five distinct filters:
\begin{enumerate}
    \item \textit{Kinematic Duration Constraints}: Assuming a circular orbit ($e=0$) and central transit ($b=0$), we calculated the maximum theoretical transit duration ($T_{\text{max}}$) using Kepler's Third Law. An observed transit duration $T_{\text{obs}}$ exceeding $T_{\text{max}}$ by $>15\%$ indicates an unphysical orbital velocity, signaling a blended binary system.
    \item \textit{Stellar Variability Check (SWEET Test)}: We applied a Lomb-Scargle periodogram to the out-of-transit data to identify dominant rotation periods.
    \item \textit{Odd-Even Depth Check}: We compared the median depths of odd and even-numbered transits to detect eclipsing binaries with primary and secondary eclipses of slightly different depths.
    \item \textit{U/V Shape Profiling}: We performed a model comparison by fitting a flat-bottomed Trapezoid (representing a planetary U-shape) and a pointed Triangle (representing a binary V-shape) to the phase-folded light curve. We computed the ratio of their Sum of Squared Residuals ($SSR_{\text{trapezoid}} / SSR_{\text{triangle}}$). A ratio $\ge 0.85$ paired with low flatness indicates a V-shaped grazing binary geometry.
\end{enumerate}

\section{Results and Discussion}
The diagnostic metrics from our vetting pipeline are compiled in Table 1.

The longer-period candidates both triggered critical physical violations. TIC 258285711 ($P = 7.5101$ d) failed the kinematic duration check ($T_{\text{obs}} = 4.08$ h vs. $T_{\text{max}} = 2.83$ h) and exhibited a phase-locked stellar rotation period of 15.00 days. TIC 319431206 ($P = 7.0008$ d) similarly exhibited an inflated duration ($T_{\text{obs}} = 4.14$ h vs. $T_{\text{max}} = 2.35$ h). Both targets are confirmed false positives.

The ultra-short-period (USP) candidates, TIC 420914536 ($P = 0.7364$ d) and TIC 432121978 ($P = 0.7123$ d), initially passed the basic kinematic duration limits. However, our advanced vetting suite successfully resolved their true natures:
\begin{itemize}
    \item \textbf{TIC 420914536}: Exhibited a severe odd-even mismatch (odd depth 0.0025 vs. even depth 0.0089). Its U/V shape fitting converged to an SSR ratio of $0.94$ with a flatness ratio of $0.01$, confirming a pointed V-shape geometry.
    \item \textbf{TIC 432121978}: Failed the odd-even check (odd depth 0.0082 vs. even depth $-0.0003$, indicating the secondary eclipse is not detected at this phase). Its U/V shape fitting yielded an SSR ratio of $0.85$ and a flatness ratio of $0.00$, confirming a grazing eclipsing binary.
\end{itemize}
Consequently, all four targets are formally retired as eclipsing binaries.

\begin{table*}[h!]
\centering
\caption{TESS Candidate Vetting, Odd-Even Checks, and U/V Shape Profiling}
\begin{tabular}{lcccccc}
\hline\hline
TIC ID & Period (d) & Obs Dur (h) & Max Dur (h) & Odd-Even Status & U/V SSR Ratio & Final Disposition \\
\hline
258285711 & 7.5101 & 4.08 & 2.83 & FAIL & 1.00 & False Positive (EB) \\
319431206 & 7.0008 & 4.14 & 2.35 & N/A & 0.98 & False Positive (EB) \\
420914536 & 0.7364 & 0.67 & 1.29 & FAIL & 0.94 & False Positive (EB) \\
432121978 & 0.7123 & 0.58 & 1.10 & FAIL & 0.85 & False Positive (EB) \\
\hline
\end{tabular}
\end{table*}

\begin{acknowledgments}
This work makes use of observations from the TESS mission and the Gaia DR3 catalog.
\end{acknowledgments}

\begin{thebibliography}{}
\bibitem[Giacalone et al.(2021)]{giacalone21} Giacalone, J., Dressing, C. D., Jensen, E. L. N., et al. 2021, AJ, 161, 24
\bibitem[Ricker et al.(2015)]{ricker15} Ricker, G. R., Winn, J. N., Vanderspek, R., et al. 2015, JATIS, 1, 014003
\end{thebibliography}

\end{document}
